\section{Introduction}
\label{sec:introduction}

Literature review is a crucial process of research. Successful literature review hinges on the ability to understand the field's landscape, present and past, and examine how the ideas are interconnected. This process can be broken down into sub-steps: retrieval (finding the right literature), analysis (understanding the connections between these papers), and summary (the synthesis of ideas and claims rooted in the source literature). The manual process is long and laborious; it begins with a few seed papers (often reviews), followed by the process of tracing citations to expand the list, and reading abstracts to decide which papers to read in full. An experienced reader needs about an hour to understand a single paper \citep{keshav2007read}, while a review on arXiv cites a median of 149 other arXiv papers and a systematic review takes 67 weeks on average \citep{huotala2024screening}.

Existing systems attempt to automate this process as search over text. Lexical and embedding retrievers are designed to retrieve literature based on term or semantic similarity \citep{robertson2009bm25,singh2022scirepeval,ostendorff2022scincl}. Deep research agents simply wrap these retrievers in a loop, where an LLM invokes queries over multiple turns \citep{he2025pasa,tongyi2025deepresearch,li2025webthinker,mirothinker2025}, treating literature review as a multi-step semantic search problem. This introduces a bottleneck reported by ScholarCatalyst \citep{kim2026scholarcatalyst}: each call queries the corpus through the same similarity retriever's top results, so repeated querying alone cannot improve over the retriever.

We argue that literature review is not simply a semantic search problem. Building on the manual process of tracing citations, systems like PaSa \citep{he2025pasa} and SciRAG \citep{ding2025scirag} attempted to leverage the citation graph in automated literature review. However, both systems incur an LLM invocation at every step for retrieval; PaSa has an LLM decide on which section's references to expand, while SciRAG uses GPT-4o to walk the citation graph. As each repeated query and citation hop invokes an LLM call, the cost of deep research systems grows quickly with the size of the query, with open agents now running up to 600 tool calls per question \citep{mirothinker2025}, at high inference cost and latency \citep{chen2026searchmore}. Yet a citation records what a field chose to read, and once the graph is built, following it needs no LLM at all.

We therefore introduce \textbf{BoundedField}, an LLM-free and citation-aware retriever that ranks papers the way a researcher reviews them, and that enhances the recall of any off-the-shelf embedding model it runs on. BoundedField takes the papers most similar to the question as a core, expands them by one citation hop, and scores each candidate as a product of bounded terms, $\prod_t (1+x_t)^{w_t}$, that measure how the core's citations reach it, along with its text match and recency. Each term lies on $[0,1]$, so a signal a paper lacks, such as citations to a paper too new to be cited, leaves its score unchanged rather than zeroing it. The exponents $w_t$ are the whole recipe: adapting BoundedField to a new task means fitting a handful of numbers on its training split, with no model trained and no LLM called, and a query takes about one second.

BoundedField runs on GraphIV, a citation graph of arXiv we computed once from arXiv's open HTML renders (3.2 million papers and 76.7 million resolved citations, growing), with each citation labeled with one of six intents and each paper with its probability of being a review. Unlike other scholarly citation graphs, which are either unfree or infrequently updated \citep{kinney2023s2,priem2022openalex,saier2023unarxive}, GraphIV is free and updates biweekly. To measure retrieval against real historical reviews, we mined three benchmarks from GraphIV's reviews, whose answer keys are reference lists that researchers wrote: SurveyMatch asks what a review cites, FrontierMatch asks for the papers the later of two reviews on the same topic cites and the earlier does not, and ConsensusMatch asks what several reviews of a topic all cite.

On these three benchmarks and five existing ones from PaSa, SurGE, and ScholarCatalyst, BoundedField led BM25, dense retrievers, a reranker, and two graph baselines at every recall cutoff but one. It more than doubled the recall of cosine search over the same embedder on all three of our benchmarks, from 0.087 to 0.216 at R@100 on SurveyMatch and from 0.257 to 0.731 on ConsensusMatch, and it beat PPR, the strongest baseline, by 51\% on SurveyMatch and 84\% on FrontierMatch. The gain comes from the graph rather than the embedder: on our three benchmarks and SurGE, BoundedField over a 475M-parameter embedder outperformed cosine search with each of seven larger embedders, up to 27B parameters \citep{zhang2025qwen3embedding,microsoft2026harrier,nvidia2026nemotronembed}. As one use of BoundedField, GraphIV-MCP gives an open-weight 27B agent six tools to search, read, and trace citations, and the agent may cite only quotes it has read (\cref{fig:teaser}). It finished all 63 report runs on a Ryzen 7 7700X and an RTX 3090 with no paid API call, at 12.7 minutes for a standard report.

Our results suggest that retrieval for literature review is better served by following what a field cites than by searching its text, and that built this way, automated literature review needs nothing more than consumer hardware. We release the graph, the benchmarks, the classifiers, and the code as a first step toward an open web of science.\footnote{\furl{https://huggingface.co/kkuroma} and \furl{https://git.kuroma.dev/kkuroma/graphiv}}
